% ============================================================
% Tufte-book compatibility shims (we migrated from tufte-book)
% ============================================================
\newcommand{\sidenote}[1]{\footnote{#1}}
\newcommand{\sidecite}[2][]{%
  \ifx\\#1\\\cite{#2}\else\cite[#1]{#2}\fi}
\newcommand{\sidenotemark}{\footnotemark}
\newcommand{\sidenotetext}[1]{\footnotetext{#1}}
\newcommand{\smallcaps}[1]{\textsc{#1}}
\newcommand{\newthought}[1]{\paragraph*{#1}}
% \sidecaption: tufte-book had a margin-caption command; map to plain \caption
\newcommand{\sidecaption}[2][]{\caption{#2}}

% ============================================================
% Encoding & basic typesetting
% ============================================================
\usepackage[utf8]{inputenc}
\usepackage[T1]{fontenc}
\usepackage{textcomp}
\usepackage{url}

\usepackage[letterpaper, margin=1in, includehead, includefoot]{geometry}

% microtype: provides \textls (letter-spacing), used in chapter title format
\usepackage{microtype}

% fmtcount: provides \NUMBERstring (spelled-out chapter numbers, e.g. "ONE")
\usepackage[english]{fmtcount}

% ============================================================
% Math (must come BEFORE cleveref)
% ============================================================
\usepackage{amsmath, amsfonts, mathtools, amsthm, amssymb}
\usepackage{physics}
\usepackage{mathrsfs}
\usepackage{cancel}

% ============================================================
% Bibliography
% ============================================================
\usepackage[
    sorting=nyt,
    style=alphabetic
]{biblatex}
\addbibresource{bibliography.bib}

% ============================================================
% Cross-referencing
% ============================================================
\usepackage{hyperref}
\hypersetup{
    colorlinks,
    linkcolor={black},
    citecolor={black},
    urlcolor={blue!80!black}
}

% ============================================================
% Tables
% ============================================================
\usepackage{multirow}
\def\block(#1,#2)#3{\multicolumn{#2}{c}{\multirow{#1}{*}{$ #3 $}}}
\usepackage{booktabs}
\usepackage{array}
\usepackage{longtable}

% caption: provides \ContinuedFloat for split figures across pages
\usepackage{caption}

% ============================================================
% Sectioning depth & ToC
% ============================================================
\usepackage{titlesec}
\usepackage{tocloft}

% Counters for the new levels
\newcounter{subsubsubsection}[subsubsection]
\newcounter{subsubsubsubsection}[subsubsubsection]
\newcounter{subsubsubsubsubsection}[subsubsubsubsection]
\newcounter{subsubsubsubsubsubsection}[subsubsubsubsubsection]
\newcounter{subsubsubsubsubsubsubsection}[subsubsubsubsubsubsection]

% Define the new sectioning levels
\titleclass{\subsubsubsection}{straight}[\subsubsection]
\titleclass{\subsubsubsubsection}{straight}[\subsubsubsection]
\titleclass{\subsubsubsubsubsection}{straight}[\subsubsubsubsection]
\titleclass{\subsubsubsubsubsubsection}{straight}[\subsubsubsubsubsection]
\titleclass{\subsubsubsubsubsubsubsection}{straight}[\subsubsubsubsubsubsection]

% Display format for the counters
\renewcommand{\thesubsubsubsection}{\thesubsubsection.\arabic{subsubsubsection}}
\renewcommand{\thesubsubsubsubsection}{\thesubsubsubsection.\arabic{subsubsubsubsection}}
\renewcommand{\thesubsubsubsubsubsection}{\thesubsubsubsubsection.\arabic{subsubsubsubsubsection}}
\renewcommand{\thesubsubsubsubsubsubsection}{\thesubsubsubsubsubsection.\arabic{subsubsubsubsubsubsection}}
\renewcommand{\thesubsubsubsubsubsubsubsection}{\thesubsubsubsubsubsubsection.\arabic{subsubsubsubsubsubsubsection}}

% Format for each level
\titleformat{\subsubsubsection}
  {\normalfont\normalsize\bfseries}{\thesubsubsubsection}{1em}{}
\titleformat{\subsubsubsubsection}
  {\normalfont\normalsize\bfseries}{\thesubsubsubsubsection}{1em}{}
\titleformat{\subsubsubsubsubsection}
  {\normalfont\normalsize\bfseries}{\thesubsubsubsubsubsection}{1em}{}
\titleformat{\subsubsubsubsubsubsection}
  {\normalfont\normalsize\bfseries}{\thesubsubsubsubsubsubsection}{1em}{}
\titleformat{\subsubsubsubsubsubsubsection}
  {\normalfont\normalsize\bfseries}{\thesubsubsubsubsubsubsubsection}{1em}{}

% Spacing for each level
\titlespacing*{\subsubsubsection}{0pt}{3.25ex plus 1ex minus .2ex}{1.5ex plus .2ex}
\titlespacing*{\subsubsubsubsection}{0pt}{3.25ex plus 1ex minus .2ex}{1.5ex plus .2ex}
\titlespacing*{\subsubsubsubsubsection}{0pt}{3.25ex plus 1ex minus .2ex}{1.5ex plus .2ex}
\titlespacing*{\subsubsubsubsubsubsection}{0pt}{3.25ex plus 1ex minus .2ex}{1.5ex plus .2ex}
\titlespacing*{\subsubsubsubsubsubsubsection}{0pt}{3.25ex plus 1ex minus .2ex}{1.5ex plus .2ex}

% TOC entries
\makeatletter
\providecommand*\l@subsubsubsection{\@dottedtocline{4}{7.0em}{4.1em}}
\providecommand*\l@subsubsubsubsection{\@dottedtocline{5}{10.0em}{5.1em}}
\providecommand*\l@subsubsubsubsubsection{\@dottedtocline{6}{13.0em}{6.1em}}
\providecommand*\l@subsubsubsubsubsubsection{\@dottedtocline{7}{16.0em}{7.1em}}
\providecommand*\l@subsubsubsubsubsubsubsection{\@dottedtocline{8}{19.0em}{8.1em}}
\makeatother

\setcounter{secnumdepth}{8}
\setcounter{tocdepth}{3}
\setlength{\cftsubsecnumwidth}{2.8em}
\setlength{\cftsubsubsecnumwidth}{3.5em}

% Cleveref must come AFTER hyperref and amsmath
\usepackage[noabbrev]{cleveref}

% Adds Bibliography to ToC
\usepackage[nottoc]{tocbibind}

% Define "normal" pagestyle as alias for "plain" (legacy tufte-book reference)
\makeatletter
\@ifundefined{ps@normal}{\let\ps@normal\ps@plain}{}
\makeatother

% ============================================================
% Graphics & color
% ============================================================
\usepackage{graphicx}
\usepackage{float}
\usepackage[usenames,dvipsnames,svgnames]{xcolor}

% ============================================================
% TikZ
% ============================================================
\usepackage{tikz}
\usepackage{tikz-cd}
\usetikzlibrary{shapes, positioning}

% ============================================================
% Theorems
% ============================================================
\usepackage{thmtools}
\usepackage{thm-restate}
\usepackage[framemethod=TikZ]{mdframed}
\mdfsetup{skipabove=1em, skipbelow=0em, innertopmargin=12pt, innerbottommargin=8pt}

\theoremstyle{definition}

\makeatletter

\declaretheoremstyle[headfont=\bfseries, bodyfont=\normalfont, mdframed={ nobreak } ]{thmgreenbox}
\declaretheoremstyle[headfont=\bfseries, bodyfont=\normalfont, mdframed={ nobreak } ]{thmredbox}
\declaretheoremstyle[headfont=\bfseries, bodyfont=\normalfont]{thmbluebox}
\declaretheoremstyle[headfont=\bfseries, bodyfont=\normalfont]{thmblueline}
\declaretheoremstyle[headfont=\bfseries, bodyfont=\normalfont, numbered=no, mdframed={ rightline=false, topline=false, bottomline=false }, qed=\qedsymbol ]{thmproofbox}
\declaretheoremstyle[headfont=\bfseries, bodyfont=\normalfont, numbered=no, mdframed={ nobreak, rightline=false, topline=false, bottomline=false } ]{thmexplanationbox}

\declaretheorem[numberwithin=chapter, style=thmgreenbox, name=Definition]{definition}
\declaretheorem[sibling=definition, style=thmredbox, name=Corollary]{corollary}
\declaretheorem[sibling=definition, style=thmredbox, name=Proposition]{prop}
\declaretheorem[sibling=definition, style=thmredbox, name=Theorem]{theorem}
\declaretheorem[sibling=definition, style=thmredbox, name=Lemma]{lemma}
\declaretheorem[sibling=definition, style=thmredbox, name=Assumption]{assumption}
\declaretheorem[sibling=definition, style=thmbluebox,  name=Example]{eg}
\declaretheorem[sibling=definition, style=thmbluebox,  name=Nonexample]{noneg}
\declaretheorem[sibling=definition, style=thmblueline, name=Remark]{remark}

\declaretheorem[numbered=no, style=thmexplanationbox, name=Proof]{explanation}
\declaretheorem[numbered=no, style=thmproofbox, name=Proof]{replacementproof}
\declaretheorem[style=thmbluebox,  numbered=no, name=Exercise]{ex}
\declaretheorem[style=thmblueline, numbered=no, name=Note]{note}

\newtheorem*{uovt}{UOVT}
\newtheorem*{notation}{Notation}
\newtheorem*{previouslyseen}{As previously seen}
\newtheorem*{problem}{Problem}
\newtheorem*{observe}{Observe}
\newtheorem*{property}{Property}
\newtheorem*{intuition}{Intuition}
\newtheorem*{rulee}{Rule}
\newtheorem*{situation}{Situation}

\declaretheoremstyle[
    headfont=\bfseries\sffamily\color{RawSienna!70!black}, bodyfont=\normalfont,
    mdframed={
        linewidth=2pt,
        rightline=false, topline=false, bottomline=false,
        linecolor=RawSienna, backgroundcolor=RawSienna!5,
    }
]{todo}
\declaretheorem[numbered=no, style=todo, name=TODO]{TODO}

\usepackage{etoolbox}
\AtEndEnvironment{vb}{\null\hfill$\diamond$}%
\AtEndEnvironment{intermezzo}{\null\hfill$\diamond$}%

\usepackage{xifthen}

\makeatother

% ============================================================
% Figure import support (Castel-style inkscape figures)
% ============================================================
\usepackage{import}
\pdfminorversion=7
\usepackage{pdfpages}
\usepackage{transparent}

\newcommand{\incfig}[2][1]{%
    \def\svgwidth{#1\textwidth}%
    \import{./figures/}{#2.pdf_tex}%
}

\newcommand{\incfigsplitting}[2][1]{%
    \def\svgwidth{#1\textwidth}%
    \import{./figures_splitting/}{#2.pdf_tex}%
}

\newcommand{\fullwidthincfig}[2][0.90]{%
    \def\svgwidth{#1\paperwidth}%
    \import{./figures/}{#2.pdf_tex}%
}

\newcommand{\minifig}[2]{%
    \def\svgwidth{#1}%
    \begingroup%
    \setbox0=\hbox{\import{./figures/}{#2.pdf_tex}}%
    \parbox{\wd0}{\box0}\endgroup%
    \hspace*{0.2cm}%
}

\pdfsuppresswarningpagegroup=1

\author{Gilles Castel}

% ============================================================
% Custom proof environment
% ============================================================
\newenvironment{myproof}[1][\proofname]{%
  \proof[\rm \bf #1]%
}{\endproof}

% ============================================================
% Algorithms & display equations
% ============================================================
\usepackage{algorithm}
\usepackage{algpseudocode}
\usepackage{breqn}

% ============================================================
% Chapter title styling (serif)
% ============================================================
\titleformat{\chapter}[display]
  {\bfseries\Large\rmfamily}
  {\filleft \textls{\MakeUppercase{\chaptertitlename}} \expandafter\textls\expandafter{\NUMBERstring{chapter}}}
  {1.5ex}
  {\titlerule
   \vspace*{1.1ex}%
   \filright}
  [\vspace*{1.5ex}%
   \titlerule]